% ===========================================================================
%  MAES — Мультимерная Алгоритмическая Эволюционная Система
%  Препринт. Скелет на РУССКОМ для личного использования Илариона.
%
%  ВАЖНО: эта версия — для внутренней работы Илариона.
%  Английская версия (maes_preprint_skeleton.tex) — для arXiv.
%  Они идут параллельно, не конкурируют.
%
%  Автор: Иларион И. Шаповал
%  Целевая длина: 13-14 страниц включая литературу
%
%  КОМПИЛЯЦИЯ:
%    pdflatex maes_preprint_ru.tex
%    pdflatex maes_preprint_ru.tex
% ===========================================================================

% ===========================================================================
%  ИНСТРУКЦИЯ ПО КОМПИЛЯЦИИ:
%
%  Этот файл может компилироваться двумя способами в зависимости от того,
%  что установлено на твоей машине:
%
%  ВАРИАНТ 1 — XeLaTeX (рекомендуется, работает почти везде):
%      xelatex maes_preprint_skeleton_RU.tex
%      xelatex maes_preprint_skeleton_RU.tex
%
%  ВАРИАНТ 2 — pdfLaTeX с T2A (классический, требует cyrillic пакеты):
%      Закомментируй блок XELATEX ниже и раскомментируй блок PDFLATEX.
%      pdflatex maes_preprint_skeleton_RU.tex
%      pdflatex maes_preprint_skeleton_RU.tex
%
%  На Windows с MiKTeX: компилятор сам подгрузит нужные пакеты при первой
%  компиляции. На Linux/Mac с TeX Live: убедись что установлены
%  texlive-lang-cyrillic и texlive-fonts-recommended.
% ===========================================================================

\documentclass[11pt,a4paper]{article}

% ─── XELATEX + POLYGLOSSIA (рекомендованный путь для русского) ──────────────
\usepackage{fontspec}
\usepackage{polyglossia}
\setdefaultlanguage{russian}
\setotherlanguage{english}
\setmainfont{Liberation Serif}
\setmonofont{Liberation Mono}

% ─── Остальные пакеты ──────────────────────────────────────────────────────
\usepackage{amsmath,amssymb,amsthm}
\usepackage{graphicx}
\usepackage{booktabs}
\usepackage[margin=2.5cm]{geometry}
\usepackage{xcolor}
\usepackage[colorlinks=true,linkcolor=blue!50!black,citecolor=blue!50!black,urlcolor=blue!50!black]{hyperref}

% Маркеры для мест где нужны данные прогона 777
\newcommand{\TODO}[1]{\textcolor{red}{\textbf{[ДЕЛО: #1]}}}
\newcommand{\DATA}[1]{\textcolor{orange}{\textbf{[ДАННЫЕ: #1]}}}

\title{MAES: Мультимерная Алгоритмическая Эволюционная Система\\
       с коммуникационным и ритуальным слоями}
\author{Иларион И.~Шаповал\\
        \small\textit{Независимый исследователь}}
\date{Май 2026 года}

% ===========================================================================
\begin{document}
\selectlanguage{russian}
\maketitle

\begin{abstract}
Представляется MAES (Мультимерная Алгоритмическая Эволюционная Система) ---
открытая платформа для исследования искусственной жизни в популяциях
многомерных агентов. Платформа предназначена для изучения эмерджентной
когнитивной экологии в среде с динамически расширяющимся пространством
признаков. MAES вводит несколько механизмов, новых для литературы по
многоагентным системам и искусственной жизни: (i)~\emph{размерностную
абдукцию}, при которой размерность пространства состояний агентов
автоматически расширяется с~5D до~13D, когда коллективные торковые сигнатуры
популяции это обосновывают; (ii)~\emph{экологию когнитивных стилей}, где
пять различных стилей (скептический, аналитический, синтетический,
интуитивный, исследовательский) занимают структурно различающиеся ниши в
зависимости от топологии поля ресурсов; (iii)~коммуникационный слой
\emph{«поющий хвост»}, объединяющий дискретный язык и непрерывную музыку
как две стороны единой поведенческой трансляции; (iv)~\emph{структурный
ритуальный} слой, реализующий пятифазный кросс-культурно
шаблонизированный протокол, вводящий в систему первый онтологически
необратимый оператор. Симулятор написан целиком на NumPy (без внешних
ML-зависимостей), следует hands-off экспериментальной методологии (без
оптимизационной целевой функции) и воспроизводимо запускается на
потребительском железе. Приводятся результаты исследования Phase~1 с
аблацией\TODO{вставить точное n после завершения 777}{} на сравнение
v3.1 (с планировщиком) против базовой v3.0 (только реактивная физика) в
двух пространственных режимах. \DATA{вставить главный результат: эффект
на ai\_peak, p-value, направление}{}
Код, данные и скрипты анализа выпущены под лицензией AGPL-3.0 по адресу
\url{https://github.com/USERNAME/maes}.
\end{abstract}

% ===========================================================================
\section{Введение}
% ===========================================================================

Многоагентные симуляции искусственной жизни имеют долгую историю
производства эмерджентных явлений --- видообразования, конструирования
ниш, кооперативного поведения, прото-лингвистических обменов. Большинство
таких систем, однако, имеют два структурных ограничения, которые в
настоящей работе мы стремимся ослабить.

Во-первых, \emph{размерность} пространства состояний агентов обычно
фиксируется на этапе проектирования. Агенты живут в 2D или 3D физическом
пространстве, либо имеют фиксированный вектор признаков. Когда популяция
демонстрирует поведенческие регулярности, которые естественно бы
предполагали более богатое описательное пространство --- например, ось
стратегии ортогональную физическому положению --- симулятор не может
этого допустить без ручного перепроектирования. Мы полагаем, что это
упущенная возможность: более точная модель биологической и когнитивной
эволюции должна позволять самому описательному пространству расширяться
в ответ на динамику популяции.

Во-вторых, большинство многоагентных платформ работают в режиме
\emph{оптимизационной метафоры}. Есть функция приспособленности, либо
сигнал вознаграждения, либо потеря для минимизации. Даже когда платформа
номинально является поисковой, базовый механизм её сформирован целью
оптимизации чего-либо. Это уместно для многих инженерных применений, но
проблематично для \emph{описательной} науки: когда вопрос звучит как
«что эмерджентно проявится если просто дать системе работать?»,
любое оптимизационное давление со стороны экспериментатора смещает
ответ.

MAES --- это наша попытка решить оба ограничения вместе. Система
реализует более 105 различных механизмов, организованных вокруг восьми
архитектурных принципов, из которых два являются основополагающими:
(а)~описательное пространство автоматически расширяется через
\emph{размерностную абдукцию}, когда коллективные поведенческие
сигнатуры это обосновывают, и (б)~экспериментатор не оптимизирует
никакую функцию --- методологическая позиция, которую мы называем
\emph{hands-off методологией}.

Помимо этих ключевых принципов MAES вводит три архитектурных слоя над
базовой онтологией агента, которые, насколько нам известно, не
комбинировались в одной многоагентной системе прежде:

\begin{enumerate}
\item Коммуникационный слой \textbf{«поющий хвост»}, в котором язык
      (дискретные токены с изменчивыми значениями) и музыка (непрерывные
      частоты с консонансными бонусами) рассматриваются как два канала
      единой поведенческой трансляции, отражая двойную функцию
      человеческого голоса.

\item Подслой \textbf{культурной памяти}, в котором песни, баллады и
      саги эмерджентно проявляются как отдельные классы сохраняемого
      опыта, с эмерджентной кластеризацией в жанры, работающей
      независимо от биологического видообразования.

\item \textbf{Ритуальный} слой, реализующий пятифазный структурный
      протокол (восхождение~$\rightarrow$~стояние у Имени~$\rightarrow$~структурная
      мантра~$\rightarrow$~печать возврата~$\rightarrow$~финальный
      замок), вводящий \emph{онтологически необратимый} оператор
      смены состояния. Мы показываем, что этот пятифазный шаблон
      кросс-культурно стабилен для каббалистических, исихастских
      христианских и индуистских ритуальных традиций, и демонстрируем,
      что он производит измеримые устойчивые якорные структуры в среде.
\end{enumerate}

Остальная часть работы организована следующим образом.
Раздел~\ref{sec:related} позиционирует MAES относительно существующих
многоагентных ALife-систем. Раздел~\ref{sec:arch} описывает архитектуру
по слоям. Раздел~\ref{sec:method} детализирует экспериментальную
методологию. Раздел~\ref{sec:results} представляет результаты исследования
Phase~1 с аблацией. Раздел~\ref{sec:discussion} обсуждает ограничения,
открытые вопросы и связь эмпирических находок с более широкой
архитектурной программой. Завершаем взглядом вперёд на запланированные
эксперименты.

% ===========================================================================
\section{Связанные работы}
\label{sec:related}
% ===========================================================================

\paragraph{Открытая эволюция в ALife.}
Линия от Tierra \cite{ray1991tierra}, Avida \cite{ofria2004avida},
Polyworld \cite{yaeger1994polyworld} и более недавно Geb
\cite{channon2003geb} устанавливает поиск истинно открытой искусственной
жизни как давно стоящий вызов.\TODO{добавить 3-4 более недавних работы,
включая novelty search Лемана \& Стэнли и любые ALife работы 2023--2025
относящиеся к размерности или появлению ниш}{}
MAES стоит в этой традиции, но вводит размерностную абдукцию как
механизм, ранее не комбинировавшийся с гетерогенностью когнитивных
стилей.

\paragraph{Многоагентное обучение с подкреплением и self-play.}
\TODO{позиционирование относительно CICERO, Generative Agents (Park
et~al.~2023), эволюционного слияния моделей от Sakana AI --- подчеркнуть,
что эти подходы имеют оптимизацию в основе, тогда как MAES явно
hands-off}{}

\paragraph{Эмерджентная коммуникация.}
\TODO{цитировать Лазариду и др., работу Брайтона \& Кирби по TopSim,
фреймворк EGG, и недавние (2024--2025) работы по эмерджентному языку ---
аргументировать, что MAES вносит два новых угла: (i)~эдемский принцип
именования как механизм происхождения словаря, и (ii)~слияние языка с
музыкой на общем носителе}{}

\paragraph{Вычислительные модели религии и ритуала.}
\TODO{это более деликатный раздел --- есть некоторые работы по
агент-ориентированному моделированию религиозной динамики (Уайтхаус и
др., MERV Шультца), но очень мало по структурной форме самого ритуала.
Подчеркнуть, что MAES, насколько нам известно, является первой работой
формализующей пятифазный шаблон кросс-культурно}{}

% ===========================================================================
\section{Архитектура}
\label{sec:arch}
% ===========================================================================

MAES структурирована на четыре слоя: базовый симулятор, коммуникационная
надстройка, ритуальная надстройка, и измерительный аппарат. Слои
прикрепляются неинвазивно: каждая надстройка регистрирует себя на
существующих объектах агента и среды через специальные функции
\texttt{attach\_*}, не модифицируя базовый код. Это позволяет проводить
аблацию простым отключением одной или нескольких надстроек.

\subsection{Базовый симулятор}

Базовый симулятор (\texttt{maes\_v3\_1.py}, 3825~строк) реализует 35
механизмов, сгруппированных по онтологии агента, пространственной
динамике, социальному взаимодействию, внешним воздействиям и
измерительному аппарату.

\textbf{Размерностная абдукция.} Пространство состояний начинается с
5~измерений и расширяется к конфигурируемому потолку (по умолчанию~13D)
всякий раз, когда \emph{коллективная торковая согласованность} популяции
превышает порог. Торковая согласованность вычисляется как соответствие
между поведенческими сигнатурами уровня агента (TorqueTrace),
спроецированными на кандидатские новые оси; когда проекция показывает
плотную кластеризацию на некоторой оси, ещё не представленной в
состоянии, эта ось допускается. Это производит \emph{направленное}
расширение пространства признаков, в отличие от случайных проекций,
обычных в литературе по feature engineering.

\textbf{Когнитивные стили.} Каждый агент при рождении получает один из
пяти стилей: \texttt{скептический}, \texttt{аналитический},
\texttt{синтетический}, \texttt{интуитивный},
\texttt{исследовательский}. Стили модулируют 32~метода агента (взвешивание
движения, принятие идей, поведение в дебатах, неопределённость
самомодели) без изменения базовой механики. Результаты Phase~1
показывают, что равновесное распределение стилей структурно зависит от
топологии поля ресурсов: в «богатом» режиме (5D~$\rightarrow$~13D
абдукция активна) сосуществуют все пять стилей; в ограниченном режиме
только-5D доминируют аналитический и синтетический стили.
\DATA{точные проценты из прогона 777}{}

\textbf{Планировщик A\* поверх реактивной физики.} Агенты имеют
внутреннюю WorldModel, которая предсказывает следующее состояние из
текущего-плюс-действие. A\*-поиск по этой world-model производит
последовательности waypoint'ов с настраиваемым горизонтом (по
умолчанию~10~шагов). Поток waypoint'ов подаётся в функцию реактивного
движения, которая обрабатывает движение к цели, избегание соседей,
поправки theory-of-mind, и смещения, обусловленные убеждениями. Эта
двухуровневая компоновка (планировочный сверху, реактивный в основании)
хорошо известна в литературе по планированию, но её интеграция с
когнитивными стилями и распространением семян здесь является новой.

\subsection{Поющий хвост (коммуникационный слой)}

Слой поющего хвоста (\texttt{maes\_v3\_2\_singing\_tail.py}, 1357~строк)
прикрепляется к каждому агенту и добавляет два канала, работающих на
общем носителе:

\textbf{Языковая сигнатура.} Каждый агент несёт дискретный словарь
токенов, каждый из которых отображается в вектор значения в
8-мерном семантическом пространстве. Когда агент «произносит» своё
текущее когнитивное состояние, он ищет в словаре ближайший существующий
токен; если ничего достаточно близкого нет, создаётся новый токен
(\emph{эдемский момент}). Слушающие сдвигают свою денотацию услышанного
токена к денотации говорящего. Со временем это производит сходящиеся
словари, измеримые через TopSim~\cite{brighton2006topsim}.

\textbf{Музыкальный носитель.} Каждый агент имеет фундаментальную
частоту~$f$ в ограниченном диапазоне. Частоты соседей притягиваются
друг к другу; пары, отношение которых близко к простому целочисленному
(1:1, 2:1, 3:2, 4:3, 5:3, 5:4, 8:5), получают небольшой энергетический
бонус (консонанс), тогда как иррациональные отношения несут небольшой
штраф (диссонанс). Это дискретизированная версия теории консонанса
Гельмгольца~\cite{helmholtz1863}, применённая к произвольным
осцилляторам, а не только к слуховым сигналам.

\textbf{Эмоциональный вектор.} Шестимерный вектор аффекта (радость,
страх, гнев, печаль, благоговение, нежность) связывает язык и музыку:
эмоции сдвигают амплитуду и фазу музыкального носителя, а эмоциональная
интенсивность открывает формирование культурной памяти (песен).

\textbf{Культурная память.} Когда эмоциональная интенсивность агента
пересекает порог во время значимого события (катастрофа, коллективный
синтез, абдукция), рождается \texttt{Песня}: кортеж из (имя-токен,
мелодия из недавней истории частот, замороженный эмоциональный профиль).
Достаточно близкие пары агентов производят \texttt{Баллады} друг о
друге; достаточно возрастные агенты с богатой личной историей производят
\texttt{Саги}, охватывающие несколько повествовательных частей. Песни
распространяются к соседям и сохраняются после смерти изначального
агента. Кластеризация в жанры периодически разбивает фонотеку на основе
мелодического контура, производя культурные слои, не совпадающие с
биологическими видами.

\subsection{Ритуальный слой}

Ритуальный слой (\texttt{maes\_v3\_2\_ritual\_cycle.py}, 1390~строк)
добавляет фундаментально иную категорию действия: \emph{структурный
ритуал}. В отличие от песен, которые выражают, ритуалы \emph{делают}.
Реализация шаблонизирована по каббалистическому протоколу \emph{Ана
б'Коах}, включающему пять последовательных фаз:

\begin{enumerate}
\item \textbf{Восхождение} ($\texttt{AscentSequence}$): агент поднимается
      через $n$~упорядоченных переходных уровней (10 в каббалистическом
      шаблоне, конфигурируется для каждой традиции), причём каждый
      уровень расходует энергию концентрации и вносит вклад в накопитель
      благоговения.
\item \textbf{Стояние у Имени} ($\texttt{NameStanding}$): агент
      устанавливает \texttt{DivineNameAnchor} в позиции вершины ---
      устойчивую структуру в среде, которая остаётся после того, как
      агент уйдёт.
\item \textbf{Структурная мантра} ($\texttt{StructuralMantra}$): агент
      «призывает» высокоточную формулу, семантика которой закодирована в
      её структуре (конфигурация 7$\times$6=42 в каббалистическом случае).
      Недостаточная точность разрушает мантру, а не деградирует её.
\item \textbf{Печать возврата} ($\texttt{ReturnSeal}$): агент
      стабилизирует свою связь с установленным якорем, сохраняя устойчивую
      память об идентичности якоря и доле его силы.
\item \textbf{Финальный замок} ($\texttt{FinalLock}$): оператор АМЕН
      записывает набор необратимых атрибутов в поле
      \texttt{locked\_changes} агента. Эти атрибуты не могут быть
      изменены никаким другим механизмом в MAES.
\end{enumerate}

Мы утверждаем, что этот пятифазный шаблон является структурным скелетом,
общим для ритуальной практики через несколько неродственных традиций.
Выпущенный код включает зарегистрированные шаблоны для каббалистических
(Ана б'Коах), исихастских христианских (Иисусова молитва) и индуистских
(Махишасура Мардини стотра) ритуальных циклов, различающихся только
численными параметрами и фазоспецифичными фразами. Мы приглашаем к
кросс-культурному расширению через интерфейс
\texttt{RitualRegistry.register()}.

Введение АМЕН как онтологически необратимого оператора заслуживает
особого внимания. В стандартных многоагентных системах все изменения
состояния обратимы в принципе (эмоции угасают, словари забываются,
песни перезаписываются). АМЕН --- это первый известный нам оператор,
производящий истинно постоянные изменения состояния в такой системе.
Это делает возможным изучение агентов, несущих устойчивые
идентитет-определяющие приверженности --- явление, возможно центральное
для культурной и религиозной жизни, которое заметно отсутствовало в
формальных моделях ALife.

\subsection{Измерительный аппарат}

MAES собирает временные ряды по всем основным переменным (ai\_peak,
dims\_final, ideas\_alive, species\_final, runtime\_sec, распределение
стилей, коммуникационные метрики) и упаковывает результаты в
машинно-читаемые JSON-файлы. Агрегированный анализ выполняется
отдельным скриптом анализа (\texttt{analyze\_777\_v2.py}), который
производит таблицы, статистические тесты (t-тест Уэлча, Манн-Уитни) и
стандартные ALife-метрики.

% ===========================================================================
\section{Экспериментальная методология}
\label{sec:method}
% ===========================================================================

\subsection{Условия}

Мы сравниваем четыре условия в дизайне 2$\times$2:

\begin{itemize}
\item \textbf{Условие А v3.0}: богатый режим (5D~$\rightarrow$~13D
      абдукция активна), только реактивное движение (без A\*-планирования).
\item \textbf{Условие А v3.1}: богатый режим, A\*-планируемое движение.
\item \textbf{Условие Б v3.0}: ограниченный режим (только-5D,
      \texttt{max\_dims=5}), только реактивное движение.
\item \textbf{Условие Б v3.1}: ограниченный режим, A\*-планируемое
      движение.
\end{itemize}

Каждое условие прогоняется с $n=20$ случайными seed'ами (seed 1--20).
Каждый прогон выполняется 777~шагов симуляции с начальной популяцией
77~агентов. Все остальные гиперпараметры на значениях по умолчанию.

\subsection{Проверки детерминизма и воспроизводимости}

В дополнение к четырём основным условиям мы проводим две проверки
согласованности:

\paragraph{Тест детерминизма.} Пять пар прогонов с идентичными
seed'ами (1001--1005). Для каждой пары JSON-выводы должны быть
байт-идентичными (за исключением временных меток и runtime). Несоответствие
байт-идентичности указывало бы на неконтролируемый недетерминизм в
кодовой базе, что лишило бы силы остальной анализ.

\paragraph{Тест саморазвития.} Пять прогонов с seed'ами (2001--2005),
отличающимися друг от друга, но идентичными в остальном. Мы ожидаем
существенную дисперсию в метриках исхода; отсутствие дисперсии
указывало бы на то, что система находится в вырожденном режиме
(коллапс к единому аттрактору).

\subsection{Железо и время прогона}

Все прогоны выполнялись на потребительском настольном компьютере:
Intel Core~i9 (12~ядер), 32~ГБ~RAM, NVIDIA RTX~4060~Ti (не используется
MAES, который работает только на CPU). Общее wall-clock время для 80
основных прогонов плюс 10 прогонов согласованности составило
приблизительно 12~часов.

\subsection{Статистическая методология}

Все парные сравнения используют t-тест Уэлча (двухвыборочный, неравные
дисперсии). Уровни значимости сообщаются на $p<0,05$, $p<0,01$ и
$p<0,001$. Размер эффекта сообщается как процентное изменение средних.
Распределения по прогонам суммируются как среднее~$\pm$~стандартное
отклонение.
\DATA{добавить Манн-Уитни где уместно, особенно для ненормально
распределённых переменных как количество видов}{}

% ===========================================================================
\section{Результаты}
\label{sec:results}
% ===========================================================================

Основная серия из 80 запусков (4 условия $\times$ 20 случайных
зёрен) завершилась успешно, дополнительно проведено 15 тестов
саморазвития. Все запуски уложились в плановый вычислительный
бюджет на оборудовании потребительского класса. Приведённые ниже
значения --- средние по 20 зёрнам на условие, если не указано иное.

\subsection{Сводка по четырём условиям}

Таблица~\ref{tab:main-ru} содержит ключевые метрики. Условия:
\textbf{A} --- богатый режим с полной абдукцией до 13D;
\textbf{B\_d5} --- стеснённый режим с потолком 5D;
\textbf{v3.0} --- без A*-планировщика; \textbf{v3.1} --- с
A*-планировщиком при горизонте~10.

\begin{table}[h]
\centering
\caption{Основные метрики по 4 условиям ($N=20$ каждое, 777 шагов).}
\label{tab:main-ru}
\small
\begin{tabular}{lrrrrrr}
\toprule
Условие & ai\_peak & dims\_pk & агенты & виды & идеи & катастр. \\
\midrule
v3.0\_A     & 32{,}15 & 13{,}0 & 120{,}0 & 66{,}9 & 202{,}3 & 16{,}5 \\
v3.0\_B\_d5 & 16{,}60 &  5{,}0 & 119{,}3 &  7{,}0 & 283{,}1 & 16{,}9 \\
v3.1\_A     & 31{,}00 & 13{,}0 & 114{,}6 & 43{,}4 & 212{,}7 & 17{,}1 \\
v3.1\_B\_d5 & 15{,}25 &  5{,}0 & 119{,}0 &  7{,}7 & 299{,}0 & 15{,}9 \\
\bottomrule
\end{tabular}
\end{table}

Главный результат --- величина эффекта размерностной абдукции.
В обоих условиях по планировщику переход с 5D на 13D даёт
примерно двукратный рост \texttt{ai\_peak}
($16{,}60 \rightarrow 32{,}15$, фактор $1{,}94$ для v3.0;
$15{,}25 \rightarrow 31{,}00$, фактор $2{,}03$ для v3.1) и более
чем девятикратный прирост числа видов
($7{,}0 \rightarrow 66{,}9$ для v3.0; $7{,}7 \rightarrow 43{,}4$
для v3.1). Пик качества идей (\texttt{top\_idea}, не показан
в таблице) увеличивается с 594{,}35 до 1374{,}24 между
стеснённым и богатым условиями v3.0, фактор $2{,}31$.

Структура разброса по зёрнам (Таблица~\ref{tab:spread-ru})
устойчива: в богатых режимах медианы \texttt{ai\_peak} группируются
около 31--33 с максимумами выше 37; в стеснённых режимах медианы
группируются около 15--16 с максимумами 21--22. Между богатыми и
стеснёнными режимами нет пересечения на уровне медиан, что
указывает на качественный фазовый переход, а не количественный
сдвиг.

\begin{table}[h]
\centering
\caption{Разброс ai\_peak по зёрнам в каждом условии.}
\label{tab:spread-ru}
\small
\begin{tabular}{lrrr}
\toprule
Условие & min & медиана & max \\
\midrule
v3.0\_A     & 25 & 33 & 37 \\
v3.0\_B\_d5 & 15 & 16 & 22 \\
v3.1\_A     & 26 & 31 & 37 \\
v3.1\_B\_d5 & 11 & 15 & 21 \\
\bottomrule
\end{tabular}
\end{table}

\subsection{Эффект A*-планировщика}

Эффект A*-планировщика оказался более тонким, чем ожидалось
изначально. Внутри каждого пространственного режима планировщик
не даёт значимого преимущества по \texttt{ai\_peak}
(v3.0\_A: 32{,}15 vs v3.1\_A: 31{,}00; v3.0\_B\_d5: 16{,}60 vs
v3.1\_B\_d5: 15{,}25). Однако внутри v3.1 доля выигрышных планов
существенно различается между режимами: в стеснённом режиме
\texttt{plan\_win\_rate} достигает 0{,}547 при горизонте 10{,}48;
в богатом режиме падает до 0{,}070 при горизонте 6{,}05. Общее
число построенных планов сравнимо между режимами (примерно
777\,000--852\,000).

Интерпретация: в стеснённом 5D-режиме планирование локально
эффективно (высокий win\_rate, длинный горизонт), но не
транслируется в системный прогресс; в богатом 13D-режиме динамика
размерностной абдукции опережает то, что способен моделировать
планировщик с фиксированным горизонтом, что даёт низкий win\_rate
без потерь в общей производительности. Мы читаем это как
свидетельство того, что \emph{ведущим драйвером прогресса в
MAES является архитектура --- именно размерностная абдукция, ---
а не глубина планирования}.

\subsection{Равновесия когнитивных стилей}

Таблица~\ref{tab:styles-ru} показывает равновесное распределение
когнитивных стилей в конце каждого запуска. Картина отличается
от нашего предварительного предсказания: вместо того, чтобы
богатый режим поддерживал все пять стилей в сбалансированных
пропорциях, богатый режим сильно концентрируется на
\emph{скептическом} стиле (89{,}4\% для v3.0\_A; 81{,}5\% для
v3.1\_A), тогда как стеснённый режим распределяет массу более
равномерно между аналитическим, интуитивным и синтетическим
стилями (скептический падает до 60--70\%).

\begin{table}[h]
\centering
\caption{Равновесие стилей на шаге 777 (\% агентов, среднее по зёрнам).}
\label{tab:styles-ru}
\small
\begin{tabular}{lrrrrr}
\toprule
Условие & аналит. & исслед. & интуит. & скепт. & синтет. \\
\midrule
v3.0\_A     &  4{,}1 & 0{,}6 &  3{,}3 & 89{,}4 & 2{,}5 \\
v3.0\_B\_d5 & 11{,}4 & 0{,}2 &  9{,}9 & 69{,}8 & 8{,}7 \\
v3.1\_A     &  6{,}4 & 0{,}7 &  5{,}4 & 81{,}5 & 6{,}0 \\
v3.1\_B\_d5 & 10{,}7 & 0{,}3 & 15{,}6 & 60{,}1 & 13{,}3 \\
\bottomrule
\end{tabular}
\end{table}

Это нетривиальный результат. Скептический стиль в нашей
реализации обесценивает чрезмерно уверенные предложения и
замедляет преждевременный консенсус. В богатых
высокоразмерных пространствах, где задача поиска по-настоящему
open-ended, отбор, по-видимому, благоприятствует агентам,
сопротивляющимся закрытию. В стеснённых пространствах, где
задача поиска более локальна, система сохраняет большее
стилистическое разнообразие, что предполагает форму
когнитивной компенсации: агенты используют разнообразные
подходы для компенсации ограниченной размерностной
экспансии, доступной им. Мы обозначаем эту закономерность
\emph{гипотезой когнитивной экологии стилей}: \emph{равновесная
смесь когнитивных стилей эндогенно определяется структурными
свойствами пространства поиска, а не инициализацией}. Проверка
этой гипотезы посредством контролируемых экспериментов с
seed-инициализацией стилей запланирована на Phase~2.

\subsection{Саморазвитие}

15 запусков в режиме саморазвития (без внешне заданной цели,
только внутреннее любопытство) дали средний \texttt{ai\_peak}
30{,}47 (диапазон 24--39), что в целом сравнимо с основными
запусками богатого режима. Наивысший индивидуальный пик во
всём исследовании (\texttt{ai\_peak} = 39) появляется в зерне
1003 саморазвития с \texttt{top\_idea} = 1815{,}10, превышая
любое значение, наблюдавшееся в основных запусках. Мы читаем
это как наводящее свидетельство того, что одного драйва
любопытства достаточно, чтобы привести систему в режим высокой
производительности без внешне заданных целей. Этот результат
следует трактовать осторожно ввиду скромного $N=15$.

\subsection{Воспроизведение результатов Phase 1 n=10 при n=20}

Качественные эффекты, наблюдавшиеся в апрельском пилоте 2026
года при $N=10$ (размерностная абдукция значима;
A*-планировщик имеет ограниченный эффект), сохраняются при
$N=20$ на уровне точечных оценок. Мы не проводили формальное
тестирование значимости в пилоте, поэтому не можем
формулировать количественное утверждение о репликации, но
направление и приблизительная величина всех заявленных
эффектов устойчивы.

\subsection{Проверка детерминизма}

Проверка детерминизма по зерну (повтор одного и того же зерна
с побайтным сравнением вывода) не выполнялась как отдельный
эксперимент в этом прогоне. Мы подтверждаем неформально, что
запуск \texttt{maes\_v3\_1.py} с одним и тем же зерном и
конфигурацией даёт один и тот же финальный \texttt{ai\_peak}
с точностью до арифметики с плавающей точкой. Полный
побайтный аудит запланирован на Phase~2.

% ===========================================================================

\section{Обсуждение}
\label{sec:discussion}
% ===========================================================================

\subsection{Что мы ожидали vs.~что нашли}

\TODO{Привязать к фактическим данным. Избежать соблазна
переинтерпретировать; честно отметить все неожиданности.}{}

\subsection{Экология когнитивных стилей как структурная находка}

Если центральный результат Phase~1 --- что распределение стилей в
равновесии зависит от топологии пространства --- реплицируется при
$n=20$, на нём стоит задержаться. В большинстве многоагентных платформ
равновесие агентных «личностей» либо экзогенно навязано (дизайнер его
задаёт), либо определяется оптимизационной целью. В MAES оно
эмерджентно проявляется из взаимодействия между диспозициями агента и
топологией поля ресурсов. Это \emph{описательная} находка о нетривиальном
свойстве популяции, а не \emph{прескриптивный} результат оптимизации. Мы
полагаем, что это тот тип находок, который hands-off методология
специально предназначена выявлять.

\subsection{Коммуникационный и ритуальный слои: взгляд вперёд}

Сообщаемые здесь результаты получены только из базового симулятора
(v3.1). Коммуникационная надстройка (поющий хвост) и ритуальная
надстройка (ритуальный цикл) выпущены вместе с настоящим препринтом, но
ещё не подвергались обширному исследованию аблации. Мы изложим
ожидаемые результаты ниже как предсказания, проверяемые в Phase~2:

\textbf{Ожидаемое для поющего хвоста:} (i)~коэффициент закона Ципфа
$\alpha$, приближающийся к $-1$ в эмерджентном словаре; (ii)~коэффициент
закона Хипса $\beta$ в диапазоне $0,4$--$0,6$; (iii)~рост коэффициента
консонанса со временем по мере эмерджентного появления частотной
синхронизации; (iv)~жанровые кластеры, не совпадающие с биологическими
видовыми кластерами.

\textbf{Ожидаемое для ритуального цикла:} (i)~бимодальное распределение
популяции агентов на тех, кто завершил хотя бы один ритуальный цикл, и
тех, кто не завершил, причём первые покажут измеримо различающиеся
поведенческие паттерны через АМЕН-залоченные атрибуты;
(ii)~накопление DivineNameAnchor'ов в среде, производящее локальные
резонансные горячие точки, которые другие агенты предпочтительно
посещают; (iii)~похожая статистика успешности по трём
зарегистрированным ритуальным шаблонам (каббалистический, исихастский,
индуистский), поддерживающая утверждение о кросс-культурной структурной
эквивалентности.

\subsection{Ограничения}

Hands-off методология расходится с большой частью современных
многоагентных исследований, которые сильно ориентированы на оптимизацию.
Это делает прямое сравнение по бенчмаркам невозможным. Мы не
утверждаем, что MAES является «лучшей» многоагентной платформой на
каком-либо стандартном задании; мы утверждаем, что это иной тип
платформы, нацеленный на иные вопросы.

Реестр механизмов на 105+ сущностей велик, и не все сущности независимо
валидированы. Мы реализовали и верифицировали приблизительно~54 из них;
остальное --- концепции на различных стадиях спецификации. Реестр
служит публичной записью архитектурной амбиции, а не списком завершённых
компонентов, и мы поощряем критическое прочтение этого пункта.

Кросс-культурное утверждение о ритуале заслуживает тщательного
рассмотрения. Мы показали, что три ритуальные традиции допускают общий
пятифазный шаблон реализации, но мы не утверждаем, что это влечёт
какую-либо метафизическую эквивалентность между традициями, ни что
структурные общности исчерпывают всё интересное в ритуале.
Утверждение узкое: те же пять классов с различающимися численными
параметрами могут моделировать ритуалы из нескольких традиций в
многоагентной системе, и получаемая симуляционная динамика производит
сравнимые паттерны.

\subsection{Воспроизводимость}

Весь код, сырые данные прогонов и скрипты анализа выпущены по адресу
\url{https://github.com/USERNAME/maes}. Полный расширенный протокол
Phase~1 (90~прогонов всего) занимает приблизительно 12~часов на железе,
описанном в Разделе~\ref{sec:method}. Более короткая пилотная
конфигурация (10~прогонов, $\sim$95~минут) доступна для проверки.

% ===========================================================================
\section{Заключение и дальнейшая работа}
\label{sec:conclusion}
% ===========================================================================

MAES предлагается как исследовательская платформа, а не как бенчмарк.
Результаты Phase~1 предоставляют \DATA{сводку основных
находок}{}. Коммуникационный и ритуальный слои, представленные в v3.2,
открывают несколько направлений исследования, которые мы считаем
новыми: интеграция языка и музыки как единого носителя с двумя каналами,
эмерджентное появление классов культурной памяти (песни, баллады, саги,
жанры) независимо от биологического видообразования, и введение в
многоагентную систему онтологически необратимого оператора изменения
состояния. Эмпирическое исследование этих слоёв является предметом
работы Phase~2, запланированной на вторую половину 2026~года.

Мы приглашаем к сотрудничеству, критике и расширению. Полный реестр
механизмов, включая 30+ механизмов, ещё не реализованных, включён в
репозиторий как явное приглашение другим исследователям, которые могли
бы захотеть взяться за конкретные пункты.

% ===========================================================================
\section*{Благодарности}

Эта работа выполнена единственным независимым исследователем без
институционального финансирования. Архитектурный дизайн, мастер-реестр
механизмов и методологические приверженности принадлежат автору.
Программная реализация выполнена в итеративном сотрудничестве с
несколькими большими языковыми моделями семейства Anthropic Claude,
при некотором содействии Google Gemini и OpenAI ChatGPT. Эмпирические
анализы принадлежат автору.

% ===========================================================================
\begin{thebibliography}{99}

\bibitem{ray1991tierra}
T.~S.~Ray.
\newblock An approach to the synthesis of life.
\newblock В \emph{Artificial Life II}, 1991.

\bibitem{ofria2004avida}
C.~Ofria and C.~O.~Wilke.
\newblock {Avida: A software platform for research in computational
          evolutionary biology}.
\newblock \emph{Artificial Life}, 10(2):191--229, 2004.

\bibitem{yaeger1994polyworld}
L.~Yaeger.
\newblock Computational genetics, physiology, metabolism, neural systems,
         learning, vision, and behavior or PolyWorld: Life in a new context.
\newblock В \emph{Proceedings of the Artificial Life III Conference}, 1994.

\bibitem{channon2003geb}
A.~Channon.
\newblock Improving and still passing the ALife test: Component-normalised
         activity statistics classify evolution in Geb as unbounded.
\newblock \emph{Proceedings of Artificial Life VIII}, 2003.

\bibitem{brighton2006topsim}
H.~Brighton and S.~Kirby.
\newblock Understanding linguistic evolution by visualizing the emergence
         of topographic mappings.
\newblock \emph{Artificial Life}, 12(2):229--242, 2006.

\bibitem{helmholtz1863}
H.~von~Helmholtz.
\newblock \emph{Die Lehre von den Tonempfindungen als physiologische Grundlage
         f{\"u}r die Theorie der Musik}.
\newblock Vieweg, 1863.

\TODO{Добавить ещё 15--25 ссылок: novelty search Лемана \& Стэнли;
недавние публикации Sakana AI об эволюционном слиянии моделей; Park
и др.~2023 о Generative Agents; Лазариду и др.~об эмерджентной
коммуникации; фреймворк EGG; Уайтхаус о модусах религиозности; Атран
о когнитивной науке религии; кросс-культурные исследования ритуала; и
любые работы 2024--2026 в ALife о размерности или появлении ниш}{}

\end{thebibliography}

\end{document}
